\chapter{Uma Inspiração Apenas}

\lettrine{H}{oje de manhã estava a conversar} com o Venerável Subbato e ele disse"-me que
nunca desenvolveu \emph{ānāpānasati}, consciência da respiração. Então eu disse
``Consegue prestar atenção a uma inspiração?'' E ele respondeu, ``Oh, sim.'' ``E
uma expiração?'' E ele respondeu, ``Sim.''

Em relação a isso, nada mais é necessário. No entanto, temos tendência de
esperar desenvolver algum tipo especial de habilidade em entrar num estado
especial. E porque não fazemos isso, pensamos então que não somos capazes de o
fazer.

Mas o caminho da vida espiritual é pelo meio da renúncia, do despojamento, do
desapego, não através do alcançar ou obter. Mesmo os
\emph{jhānas}\footnote{\emph{Jhānas}: são estados sublimados de
  mente"-consciência experienciados por meio da absorção meditativa.} são
renúncias, mais do que conquistas. Se renunciarmos mais e mais, bem como
abrirmos mão mais e mais, então os estados \emph{jhānicos} serão naturais.

A atitude é extremamente importante. Para se praticar \emph{ānāpānasati}, a
pessoa leva a atenção até uma inspiração, estando consciente do princípio ao
fim. Uma inspiração, é isso; e então procede"-se da mesma forma para a expiração.
Essa é a realização perfeita de \emph{ānāpānasati}. A consciência precisamente
só disto é o resultado da concentração da mente pela atenção sustentada na
respiração. Do início ao fim da inspiração, do início ao fim da expiração. A
atitude é sempre a de abrir mão, não nos apegando a quaisquer ideias ou
sentimentos que surjam a partir daí, de modo a estarmos sempre renovados com a
próxima inspiração, a próxima expiração, completamente tal como é. Não estamos a
transportar nada. Portanto, é uma forma de abandonar, de largar, em vez de
atingir e alcançar.

Os perigos na prática da meditação são os hábitos de cobiçarmos coisas,
cobiçarmos estados; então o conceito que é mais útil é o conceito de abrir mão,
ao invés de atingir e alcançar. Se dissermos hoje que ontem tivemos uma super
meditação, absolutamente fantástica, exactamente como sempre sonhámos, e hoje
tentamos obter a mesma experiência maravilhosa de ontem, ficamos mais inquietos
e agitados que nunca. Ora porque é que isso acontece? Por é que não podemos
alcançar o que queremos? É porque estamos a tentar alcançar algo de que nos
lembramos, em vez de realmente trabalharmos com a realidade como as coisas são,
como estão a ocorrer no agora. Portanto, a via correcta é uma via de
consciência, de olhar para a realidade como é agora, em vez de lembrarmo"-nos de
ontem e tentarmos alcançar aquele estado outra vez.

No primeiro ano em que meditei, não tive um professor. Estive num pequeno
\emph{kuṭi}\footnote{\emph{Kuṭi}: uma cabana de madeira sem mobília muito
  simples que serve de abrigo para um monge ou monja budista.} em Nong Khai
durante cerca de dez meses e tive todos os tipos de visões brilhantes. Ficar dez
meses sozinho, sem ter que falar, sem ter que ir a lugar algum, tudo se acalmou
após vários meses, e então pensei que era uma pessoa totalmente iluminada, um
\emph{arahant}. Eu tinha a certeza disso. Mais tarde descobri que não era.

Lembro"-me que naquele ano passámos por uma escassez de comida em Nong Khai e não
tivemos muito para comer. Eu sofri de desnutrição, e na altura pensei, ``Talvez
a desnutrição seja a resposta. Se eu simplesmente me abster de comer\ldots{}'' Eu
lembro"-me de estar tão fraco com desnutrição em Nong Khai, que as minhas orelhas
começaram a rachar. Ao acordar, eu tinha que forçar as minhas pálpebras para
abri"-las; estavam coladas com o que sai das pálpebras quando não estamos muito
bem.

Então, um dia, um monge canadense trouxe"-me três latas de leite enlatado. Na
Ásia eles têm leite açucarado em lata e é muito saboroso. Ele trouxe"-me também
um pouco de café instantâneo e uma garrafa de água quente. Então fiz uma bebida:
coloquei um pouco de café, pus um pouco desse leite, verti água quente e comecei
a beber. Eu simplesmente fiquei louco. Foi tão absolutamente delicioso, a
primeira vez que tive algo doce ou estimulante durante semanas. E estando eu
desnutrido e num estado apático de cansaço extremamente embotado, isto era como
gasolina de elevada octanagem -- zás! Engoli de imediato -- eu não conseguia
parar -- e consegui consumir todas as três latas de leite e uma boa parte
daquele café. E a minha mente verdadeiramente voou para o espaço sideral, ou
assim parecia, e pensei, ``Talvez seja esse o segredo. Se eu conseguir que
alguém me compre leite enlatado.''

Quando fui para Wat Pah Pong no ano seguinte, continuava a pensar, ``Oh, tive
todas aquelas experiências maravilhosas em Nong Khai. Tive todo aquele leque
maravilhoso de visões, aquelas experiências fantásticas de flutuação e visões
brilhantes, e parecia que eu compreendia tudo. E até pensaste que eras um
arahant.'' Em Wat Pah Pong, naquele primeiro ano, não tive muito. Continuei
simplesmente a tentar fazer tudo o que tinha feito em Nong Khai para conseguir
essas coisas. Mas depois de algum tempo, mesmo com fortes chávenas de café, já
não funcionou. Eu parecia já não estar a conseguir aquelas alegrias, aquelas
alturas fantásticas e visões brilhantes que tive no primeiro ano. Então, depois
do primeiro Vassa\footnote{Vassa: o tradicional retiro das chuvas de três meses
  realizado a cada ano em mosteiros budistas. Geralmente é um momento de maior
  atenção às questões do treino e instrução espiritual.} em Wat Pah Pong,
pensei: ``Este lugar não é para mim. Acho que vou tentar repetir o que aconteceu
em Nong Khai.'' Deixei Ajahn Chah e fui viver na montanha Pupek na província de
Sakorn Nakorn.

Lá, finalmente, eu estava num local idílico. No entanto, para a ronda da
mendicância, era preciso sair antes do amanhecer e descer a montanha, que era
uma subida e tanto, e esperar a chegada dos aldeões. Eles traziam"-me comida, e
então tinha que subir todo o caminho de volta e comer essa comida antes do
meio"-dia. Isso foi um grande problema.

Eu estava com um outro monge, um monge tailandês, e pensei, ``Ele é realmente
muito bom'' e fiquei bastante impressionado com ele. Mas quando estávamos nesta
montanha, ele queria que eu lhe ensinasse inglês -- então fiquei mesmo
enraivecido com ele e tive vontade de o matar.

Foi numa área onde havia muitos terroristas e comunistas, no nordeste da
Tailândia. Às vezes havia helicópteros a voar sobre as nossas cabeças
vigiando"-nos. Uma vez, vieram e levaram"-me para a cidade da província, admitindo
que eu poderia ser um espião comunista.

Então adoeci violentamente, tão mal, que tiveram que me carregar montanha
abaixo. Eu estava enfiado num lugar miserável, perto de um reservatório sob
telhado de zinco, na estação quente, com insectos a zumbir dentro e fora das
minhas orelhas e orifícios. Com comida horrível. Quase morri, pensando bem.
Quase que não sobrevivi.

Mas foi durante esse período naquele alpendre de telhado de zinco que ocorreu
uma verdadeira mudança. Eu estava realmente desesperado, doente, fraco e
totalmente deprimido, e a minha mente caia nesses reinos infernais, com o
terrível calor e desconforto. Eu senti"-me como se estivesse sendo cozinhado; foi
tipo tortura.

Então uma mudança sobreveio. De repente, simplesmente parei a minha mente;
recusei"-me a ser apanhado nessa negatividade e comecei a praticar
\emph{ānāpānasati}. Usei a respiração para concentrar a minha mente e as coisas
mudaram extremamente rápido. Depois disso, recuperei a minha saúde e chegou a
hora de entrar no Vassa seguinte -- então regressei -- tinha prometido a Ajahn
Chah que voltaria a Wat Pah Pong para passar o Vassa -- e as minhas vestes
estavam todas esfarrapadas, rasgadas e remendadas. Eu estava com um aspecto
terrível. Quando Ajahn Chah me viu, simplesmente explodiu à gargalhada. E fiquei
tão feliz por voltar depois daquilo tudo!

Eu tinha tentado praticar e o que tinha desejado eram as memórias dessas
compreensões lúcidas. Tinha esquecido o que as compreensões lúcidas realmente
eram. Eu estava mesmo muito apegado à ideia de trabalhar nalgum tipo de caminho
ascético como no primeiro ano, quando o ascetismo realmente funcionou. Naquela
época, estar desnutrido e sozinho pareceu ter"-me provido de compreensão lúcida,
de modo a que, nos anos seguintes, continuei a tentar criar as condições em que
pudesse ter essas fantásticas compreensões lúcidas.

Mas os dois ou três anos seguintes pareceram anos de apenas passar o tempo. Nada
parecia acontecer. Estive seis meses nesta montanha antes de regressar a Wat Pah
Pong, apenas decidindo ficar e seguir as compreensões lúcidas que tive. Uma das
compreensões lúcidas do primeiro ano foi que eu deveria encontrar um professor e
que deveria aprender a viver sob uma disciplina imposta por esse professor.
Então fiz isso. Percebi que Ajahn Chah era um bom professor e que tinha um bom
nível de disciplina monástica, e portanto permaneci com ele. Essas compreensões
lúcidas que tive estavam certas, mas eu tinha"-me apegado à memória.

As pessoas ficam muito apegadas a todas essas coisas especiais, como retiros e
cursos de meditação onde tudo está sob controlo, e tudo é organizado e há
silêncio total. Então, ainda que tenhamos compreensão lúcida, a reflexão nem
sempre está presente porque presumimos que, para ter essas compreensões lúcidas,
precisamos dessas condições.

Na realidade, a compreensão lúcida é cada vez mais uma questão de viver com
perspicácia. Não é apenas às vezes ter uma compreensão lúcida, mas cada vez
mais, à medida que reflectimos no Dhamma, tudo então é profundo. Observamos a
vida de forma perspicaz enquanto ocorre connosco. Mal achamos que precisamos de
condições especiais para isso, sem estarmos cientes disso, criamos todos os
tipos de complexidades sobre a nossa prática.

Então desenvolvi o abrir mão: o não me preocupar em alcançar ou realizar nada.
Decidi tornar possíveis pequenas conquistas ao aprender a ser um pouco mais
paciente, um pouco mais humilde e um pouco mais generoso. Decidi desenvolver
isto, em vez de sair do meu caminho para controlar e manipular o ambiente, com a
intenção de me preparar na esperança de ficar extasiado. Tornou"-se óbvio, por
meio de reflexão, que o apego às compreensões lúcidas era o problema. Estas eram
compreensões lúcidas válidas, mas havia apego à memória.

Logo surgiu"-me a compreensão lúcida de abrir mão de todas as compreensões
lúcidas. Não nos apegamos a elas. Simplesmente continuamos a deixar de lado
todas as compreensões lúcidas que temos, porque caso contrário elas tornam"-se
memórias, e então as memórias são condições da mente e se nos apegamos a elas,
só nos podem levar ao desespero.

Em cada momento, é como é. Com \emph{ānāpānasati}, neste momento, uma
inspiração, é desta forma. Não é como foi a inspiração de ontem. Não estamos a
pensar na inspiração ou expiração de ontem enquanto fazemos a de agora. Estamos
completamente com esta, tal como é; então nós estabelecemos isso. A capacidade
reflexiva é baseada em estabelecer a nossa consciência na realidade como é
agora, ao invés de ter alguma ideia do que gostaríamos de obter e, em seguida,
tentar obter isso aqui e agora. Tentar obter a sensação de felicidade de ontem
no aqui e agora significa que não estamos cientes da realidade do agora. Não
estamos com ela. Mesmo com \emph{ānāpānasati}, se praticarmos com a esperança de obter
o resultado que tivemos ontem, isso tornará impossível que esse resultado ocorra
novamente alguma vez.

No inverno passado, o Venerável Vipassi estava a meditar no templo e alguém
estava a fazer ruídos bastante perturbadores. Ao conversar com o Venerável
Vipassi sobre isso, fiquei bastante impressionado, porque ele disse primeiro que
se sentia aborrecido e depois decidiu que os ruídos fariam parte da prática.
Então, ele abriu a sua mente para a sala de meditação com tudo incluído -- os
ruídos, o silêncio, a coisa toda. Isso é sabedoria, não é? Se o barulho é algo
que podemos parar -- como uma porta a bater ao vento -- vamos e fechamos a
porta. Se houver algo sobre o qual temos controlo, isso podemos fazer.

Mas nós não temos controlo sobre grande parte da vida. Não temos o direito de
pedir a tudo para ficar em silêncio para a ``minha'' meditação. Quando há
reflexão, ao contrário de termos uma pequena mente que tem que ter silêncio
total e condições especiais, temos uma grande mente que pode conter tudo isso:
os ruídos, as perturbações, o silêncio, a felicidade, a inquietação, a dor. A
mente é toda abrangente, em vez de se especializar num certo refinamento da
consciência. Assim desenvolvemos flexibilidade porque podemos concentrar a nossa
mente.

É aqui que a sabedoria é necessária para o verdadeiro desenvolvimento. É por
meio da sabedoria que ele se desenvolve, não através da força de vontade ou do
controlo ou manipulação das condições ambientais. Ou de nos vermos livres das
coisas que não queremos e preparar"-nos para que possamos seguir esse desejo de
atingir e alcançar. O desejo é insidioso. Quando estamos cientes de que a nossa
intenção é atingir algum estado, isso é um desejo, não é? Então abrimos mão
disso. Se estamos sentados aqui, mesmo com o desejo de atingir o primeiro
\emph{jhāna}, reconhecemos que esse desejo será exactamente o que vai impedir a
realização. Então, abandonamos o desejo, o que não significa não fazer
\emph{ānāpānasati}, mas mudar a atitude em relação a isso.

Então, o que é que podemos fazer agora? Desenvolver consciência de uma
inspiração. A maioria de nós pode fazer isso; a maioria dos seres humanos tem
concentração suficiente para se concentrar do início ao fim de uma inspiração.
Mas mesmo que a nossa concentração seja tão fraca que nem sequer consigamos
chegar ao fim, não há problema. Pelo menos podemos chegar ao meio, talvez. Isso
é melhor do que se tivéssemos desistido totalmente ou nunca tivéssemos tentado,
não é? Porque pelo menos estamos a aquietar a mente por um segundo, e esse é o
começo: aprender a acalmar e a concentrar a mente em torno de uma coisa, como a
respiração, e mantê"-la simplesmente durante uma inalação; se não, então de meia
inalação, ou um quarto, ou o que for. Pelo menos começámos e devemos tentar
desenvolver uma mente que fica feliz por simplesmente ser capaz de fazer esse
tanto, em vez de criticarmos porque é que não atingimos o primeiro \emph{jhāna},
ou o quarto.

Se a meditação se tornar outra coisa que devemos fazer e nos sentimos culpados
por não vivermos à altura das nossas resoluções, então começamos a forçar sem
termos consciência do que estamos a fazer. Então a vida fica deveras triste e
deprimente. Mas se estivermos a colocar esse tipo de atenção inteligente na
nossa vida diária, vamos descobrir que muito da vida diária é bastante
agradável, mas que podemos não notar isso se estivermos presos nas nossas
compulsões e obsessões. Se agirmos com compulsão, isso torna"-se um fardo, uma
opressão. Depois, arrastamo"-nos por aqui e por acolá, cumprindo o que temos que
fazer de uma forma descuidada e negativa. Mas tendo a possibilidade de estarmos
no campo -- as árvores, a natureza\ldots{} no entanto temos este tempo para um
retiro -- podemos sentar e caminhar, e não temos muito que fazer. Os cânticos
matinais, os cânticos nocturnos, podem ser extremamente agradáveis, quando
estamos abertos a isso. As pessoas estão a oferecer a comida. A refeição é algo
deveras bonito.

As pessoas estão a comer conscientemente e em silêncio. Quando estamos a fazer
isso por hábito e compulsão, então torna"-se aborrecido. E uma série de coisas
que normalmente são bastante agradáveis, deixam de o ser. Não conseguimos
apreciá"-las quando funcionamos a partir da compulsividade, inconsciência e
ambição. Estes são os tipos de forças motrizes que destroem a alegria e o
encanto nas nossas vidas.

Manter a nossa atenção na respiração, de facto desenvolve a consciência. Mas
quando nos perdemos no pensamento ou na inquietação, isso também faz parte e não
é problema. Não nos podemos pressionar demasiado. Não sejamos tiranos nem nos
castiguemos com um chicote nem nos forcemos de forma maldosa. Possamos liderar,
orientar e treinar; liderando para a frente, orientemo"-nos em vez de nos
contorcermos ou forçarmos. Nibbāna é uma compreensão subtil de não possuir. Não
é possível conduzirmo"-nos ao Nibbāna. Essa é a via mais certa para nunca o
realizarmos. É aqui e agora, portanto se nos estamos a dirigir para o Nibbāna,
estamos sempre a distanciar"-nos dele, falhando a sua realização.

É muito duro, por vezes, queimar apegos na nossa mente. A Vida Sagrada é um
holocausto, uma queima total, uma queima do eu, da ignorância. Um diamante é um
símbolo da pureza que vem do holocausto; algo que ao atravessar tais incêndios,
tudo o que restou foi a pureza. E então, é por isso que nesta nossa vida aqui,
tem que haver esta vontade de queimar os preconceitos egocêntricos, as opiniões,
os desejos, a inquietação, a ganância, tudo isso, toda essa totalidade, de modo
que não haja nada para além da pureza remanescente. Então, quando há pureza, não
há ninguém, nada, há isso, aquilo ``que é''.

E abramos mão disso. Mais e mais o caminho é simplesmente o saber estar aqui e
agora, saber estar com a realidade assim como as coisas são. Não há para onde
ir, nada para fazer, nada para se tornar, nada para nos desfazermos. Por causa
do holocausto, não há ignorância remanescente; existe pureza, clareza e
inteligência.
